\documentclass[11pt]{amsart}
\usepackage[T1]{fontenc}
\usepackage[utf8]{inputenc}
\usepackage[french]{babel}
\usepackage{amssymb}
\usepackage[margin=1.05in]{geometry}
\usepackage[hidelinks]{hyperref}
\newtheorem{proposition}{Proposition}
\newtheorem{lemma}[proposition]{Lemme}
\theoremstyle{remark}
\newtheorem{remark}[proposition]{Remarque}
\DeclareMathOperator{\im}{im}
\newcommand{\PP}{\mathbf P}
\newcommand{\OO}{\mathcal O}
\title{Des générateurs explicites pour la cohomologie nécessaire au relèvement}
\date{26 octobre 2026}
\begin{document}
\maketitle

\noindent\textbf{Statut.}
Cette note complète l'argument de \texttt{finite-lifting.tex}.
Elle construit les générateurs de cohomologie nécessaires à partir
d'une liste finie de générateurs d'un idéal homogène.
Dans le cas $\mathbf Q[a,b]$, cet idéal s'obtient par une élimination
polynomiale, couverte par les méthodes constructives de
Coquand--Persson~\cite{CP}.
Ce sont des preuves et des algorithmes informels; le développement
de Čech, l'élimination et leur application ne sont pas encore
formalisés dans Unit.

\section{Un cas où une présentation suffit}

Soient $D$ un anneau commutatif quelconque, $S=D[X,Y,Z]$, gradué
par le degré total en $X,Y,Z$, et
\[
 J=(g_1,\ldots,g_r),\qquad \deg(g_j)=d_j\geq0,\qquad N=S/J.
\]
Les générateurs sont homogènes; un générateur nul peut être supprimé
ou simplement produire des classes nulles. Tous les calculs de
cohomologie ci-dessous sont ceux du complexe de Čech des trois cartes
standard sur $\PP^2_D$.

\begin{proposition}\label{prop:finite}
Il existe une surjection explicite
\[
 \bigoplus_{j=1}^r D^{\binom{d_j}{2}}
 \longrightarrow H^1(\PP^2_D,\widetilde N(-1)),
 \qquad \binom d2=0\text{ pour }d<2.
\]
Cette construction ne suppose ni $D$ noethérien, ni $J$ engendré par
une suite régulière, ni une résolution libre de $N$.
\end{proposition}

La raison cohomologique est courte. La suite
\[
 0\longrightarrow\widetilde J(-1)\longrightarrow\OO(-1)
 \longrightarrow\widetilde N(-1)\longrightarrow0
\]
donne $H^1(\widetilde N(-1))\simeq H^2(\widetilde J(-1))$, car
$H^1(\OO(-1))=H^2(\OO(-1))=0$.
La surjection $\bigoplus_j\OO(-d_j-1)\to\widetilde J(-1)$ donne
une surjection sur $H^2$, puisque le complexe de Čech n'a pas de
terme en degré $3$. Enfin
\[
 H^2(\OO(-d-1))\simeq D^{\binom d2}.
\]
La section suivante donne les représentants et une preuve constructive
de la surjectivité, sans utiliser un théorème général de finitude.

\section{Les cocycles, monôme par monôme}

Écrivons $C^1=(C_{XY},C_{XZ},C_{YZ})$ et
\[
 \delta(b)=b_{YZ}-b_{XZ}+b_{XY},\qquad
 (\delta c)_{ij}=c_j-c_i.
\]
Sur la triple intersection, les éléments sont des polynômes de Laurent
en $X,Y,Z$ à coefficients dans $D$, de degré total prescrit.

Pour un tel polynôme $q$, soit $\pi(q)$ la somme des monômes dont
les trois exposants sont strictement négatifs. On définit $\sigma(q)$
en traitant chaque autre monôme $m=X^{e_0}Y^{e_1}Z^{e_2}$ : choisir
le premier indice $i$ avec $e_i\geq0$, et placer $(-1)^i m$ dans la
composante double qui omet $i$. Alors
\[
 q=\pi(q)+\delta\sigma(q).
 \tag{1}
\]
Cette règle n'utilise que la comparaison d'entiers; elle ne teste pas
si un coefficient de $D$ est nul.

En degré total $-1$, $\pi(q)=0$, donc $\delta\sigma(q)=q$.
De plus, toute $1$-cochaîne fermée, de n'importe quel degré, est le
bord d'une $0$-cochaîne. Voici une règle explicite $\tau$, appliquée
à chaque monôme. Si $b=(b_{XY},b_{XZ},b_{YZ})$ est son triplet de
coefficients et $E$ l'ensemble des indices d'exposant négatif, poser
\[
\begin{array}{c|c}
 E&\tau(b)=(c_X,c_Y,c_Z)\\\hline
 \varnothing&(0,b_{XY},b_{XZ})\\
 \{X\}&(-b_{XY},0,0)\\
 \{Y\}&(0,b_{XY},0)\\
 \{Z\}&(0,0,b_{XZ})\\
 |E|\geq2&(0,0,0).
\end{array}
\]
Pour $|E|=2$, une seule composante double est possible, et la
fermeture force son coefficient à être nul; pour $|E|=3$, il n'y
en a aucune. Dans chaque cas, $\delta\tau(b)=b$ dès que $\delta b=0$.
Ce sont les contractions usuelles de Čech~\cite{proj}.

Pour $g_j$ et chaque triplet
\[
 \alpha,\beta,\gamma\geq1,\qquad
 \alpha+\beta+\gamma=d_j+1,
\]
posons
\[
 \ell_{j,\alpha\beta\gamma}=X^{-\alpha}Y^{-\beta}Z^{-\gamma},
 \qquad
 b_{j,\alpha\beta\gamma}=\sigma(g_j\ell_{j,\alpha\beta\gamma}).
\]
Le produit $g_j\ell_{j,\alpha\beta\gamma}$ est de degré $-1$,
donc $\delta b_{j,\alpha\beta\gamma}=g_j\ell_{j,\alpha\beta\gamma}$.
Après réduction modulo $J$, $b_{j,\alpha\beta\gamma}$ est donc
un $1$-cocycle. Il y a exactement $\binom{d_j}{2}$ tels triplets.
Leurs classes sont les générateurs annoncés.

\begin{proof}[Preuve explicite de la surjectivité]
Soit $b$ un $1$-cocycle à valeurs dans $\widetilde N(-1)$.
Relevons ses trois composantes en une cochaîne $\widetilde b$ de
$\OO(-1)$. La fermeture modulo $J$ donne une écriture finie
\[
 \delta\widetilde b=\sum_j g_j q_j,
 \qquad \deg(q_j)=-d_j-1,
 \tag{2}
\]
dans l'anneau de Laurent de la triple intersection.
Cette écriture est un témoin d'appartenance à l'idéal localisé.
Par (1), écrire $q_j=\pi(q_j)+\delta u_j$, où $u_j=\sigma(q_j)$.
Chaque $\pi(q_j)$ est combinaison $D$-linéaire des $\ell_{j,\alpha\beta\gamma}$.
Posons
\[
 b_* = \sum_j \sigma(g_j\pi(q_j)).
\]
Les équations (1)--(2) donnent
\[
 \delta\left(\widetilde b-\sum_jg_j u_j-b_*\right)=0.
\]
La contraction en degré $1$ de $\OO(-1)$ produit $c$ tel que
\[
 \widetilde b-\sum_jg_j u_j-b_*=\delta c.
\]
Modulo $J$, cela prouve que $b$ et $b_*$ représentent la même classe.
Toutes les décompositions utilisées sont finies.
\end{proof}

\section{La présentation du module de Rees}

Revenons à $A=\mathbf Q[a,b]$, $I=(a,b)$,
$S_0=A[X,Y,Z]$, et aux homogénéisés $P,Q$ de $p,q$.
Posons $B_h=S_0/(P,Q)$ et $W=\operatorname{Proj}(B_h)$.
Le module de Rees de $B_h$ est ici une algèbre cyclique :
\[
 \mathcal N=\bigoplus_{n\geq0}I^nB_hT^n
 =\im\bigl(A[U,V,X,Y,Z]\longrightarrow B_h[T]\bigr),
 \quad U\mapsto aT,\quad V\mapsto bT.
\]
Avec $D=A[U,V]$, sa présentation est $\mathcal N\simeq D[X,Y,Z]/J$,
où
\[
 J=(P,Q,U-aT,V-bT)\,\mathbf Q[a,b,U,V,X,Y,Z,T]
       \ \cap\ \mathbf Q[a,b,U,V,X,Y,Z].
 \tag{3}
\]
Cette égalité résulte du quotient par le graphe de la substitution.
Il faut utiliser tout l'idéal (3); les seuls $P,Q,bU-aV$ ne sont
pas en général une présentation suffisante.

Les deux degrés à conserver sont
\[
\begin{array}{c|rrrrrrrr}
 &a&b&U&V&T&X&Y&Z\\\hline
 \text{projectif}&0&0&0&0&0&1&1&1\\
 \text{Rees}&0&0&1&1&1&0&0&0
\end{array}
\]
Les quatre équations de l'idéal avant élimination sont bihomogènes.
On calcule une base de Gröbner avec un ordre d'élimination pour $T$,
par exemple lexicographique avec $T$ en première position.
Les opérations de réduction préservent les deux homogénéités.
Les éléments de la base qui ne contiennent pas $T$ engendrent $J$.
L'argument d'élimination est immédiat sur la division : le monôme
dominant d'un polynôme sans $T$ ne peut être divisible par un monôme
contenant $T$; avec cet ordre, un élément de la base dont le monôme
dominant ne contient pas $T$ n'en contient nulle part.

L'existence constructive d'une base de Gröbner sur un corps discret
est le théorème 11 de Coquand--Persson~\cite{CP}.
On obtient ainsi une liste finie bihomogène $g_j$, avec degrés
projectifs $d_j$ et degrés de Rees $e_j$.

Sur chaque carte et intersection de cartes, localisation en $X,Y,Z$
et extraction du degré projectif $-1$ donnent
\[
 \check C^\bullet(\mathcal N(-1))
  =\bigoplus_{n\geq0}\check C^\bullet(W,I^n\OO_W(-1))T^n.
 \tag{4}
\]
La localisation commute aux sommes directes et aux images de
multiplication; les trois cartes donnent des produits finis.
La cohomologie de cette somme directe se calcule degré par degré,
puisque tous les éléments ont un support fini.

La proposition~\ref{prop:finite} appliquée à (3)--(4) fournit donc
des générateurs explicites de
\[
 \mathcal M=\bigoplus_{n\geq0}H^1(W,I^n\OO_W(-1))T^n.
\]
Ils ont les degrés de Rees $e_j$. Le module est engendré sur $D$
par au plus $\sum_j\binom{d_j}{2}$ éléments. Comme l'action de $D$
se factorise par $A[U,V]/(bU-aV)=\mathcal R$, la même liste engendre
$\mathcal M$ sur $\mathcal R$.

\section{La finitude des sections globales pour deux équations}

Pour terminer le passage aux éléments entiers, il faut aussi produire
des générateurs de $C=H^0(W,\OO_W)$ sur $A$.
Ici encore une résolution générale n'est pas nécessaire.
Dans l'anneau factoriel $S_0=\mathbf Q[a,b,X,Y,Z]$, calculons un
pgcd homogène $D_0$ de $P,Q$. Posons
\[
 d=\deg P,\quad e=\deg Q,\quad r=\deg D_0,
 \quad P'=P/D_0,\quad Q'=Q/D_0.
\]
La coprimalité de $P',Q'$ donne la suite exacte graduée
\[
 0\longrightarrow S_0(-d-e+r)
 \xrightarrow{\binom{-Q'}{P'}} S_0(-d)\oplus S_0(-e)
 \xrightarrow{(P\ Q)} S_0
 \longrightarrow B_h\longrightarrow0.
 \tag{5}
\]
En effet $Pu+Qv=0$ implique $P'u+Q'v=0$, puis
$u=-Q'w$, $v=P'w$. Il n'est pas supposé que $P,Q$ soient une suite
régulière.

Soit $L$ le noyau de l'application entre modules libres finis sur $A$
\[
 H^2(\OO(-d-e+r))
   \xrightarrow{\binom{-Q'}{P'}}
 H^2(\OO(-d))\oplus H^2(\OO(-e)).
 \tag{6}
\]
Les entrées de cette matrice se calculent en multipliant les monômes
de Laurent puis en gardant les exposants tous négatifs.
Puisque $d,e\geq1$ et $d+e-r\geq1$, (5) et le calcul de Čech donnent
\[
 0\longrightarrow A\longrightarrow C\longrightarrow L\longrightarrow0.
 \tag{7}
\]
Des générateurs du noyau de la matrice (6) se calculent sur
$A=\mathbf Q[a,b]$ par élimination. Leurs relèvements dans $C$
s'obtiennent explicitement : représenter $t\in L$ par une somme
de monômes de Laurent tous négatifs de degré $-d-e+r$. Alors
\[
 \pi(-Q't)=\pi(P't)=0,\qquad
 u_1=\sigma(-Q't),\quad u_2=\sigma(P't).
\]
On a $\delta(Pu_1+Qu_2)=0$, donc
\[
 \overline{\tau(Pu_1+Qu_2)}\in C
\]
est un relèvement de $t$ dans (7). Avec $1$, ces relèvements
engendrent $C$. Pour l'exemple quadratique d'ACL,
$d=e=2,r=0$ et (6) est $A^3\to0$; (7) retrouve le rang $4$.

On peut réduire le calcul du noyau d'une matrice $T:A^m\to A^s$
au seul calcul de bases de Gröbner d'idéaux : introduire $w_1,\ldots,w_m$
et $z_1,\ldots,z_s$, puis éliminer les $z_j$ de
\[
 (z_jz_k,\ w_i-\sum_j T_{ji}z_j)\subseteq A[\mathbf w,\mathbf z].
\]
Le noyau de la substitution vers $A[\mathbf z]/(\mathbf z)^2$
est homogène pour $\deg w_i=\deg z_j=1$ et n'a pas d'élément
non nul de degré zéro. Les générateurs de degré $1$ de l'idéal
d'élimination sont exactement une famille génératrice de $\ker T$,
en identifiant $\sum_i c_iw_i$ à $(c_i)$.

\section{Conséquences pour la formalisation}

La preuve complète envisagée pour $p\equiv x$, $q\equiv y\pmod{(a,b)}$
peut donc suivre les étapes suivantes.
\begin{enumerate}
\item Construire l'idéal de Rees (3) par élimination.
\item Construire les cocycles de la proposition~\ref{prop:finite}
et leur propriété de génération par les contractions monomiales.
\item Appliquer la relation finie entre obstructions et la correction
de \texttt{finite-lifting.tex} à la cochaîne fixe $(0,0,Z^{-1})$
pour obtenir $h$. L'identité $x,y\in(a,b)B$ fournit ses relèvements
à tous les ordres par les mêmes représentants.
\item Construire $C$ par (5)--(7), puis appliquer le déterminant
aux multiplications par $Zh,Xh,Yh$ pour obtenir les équations unitaires.
\end{enumerate}
Les étapes sont des constructions algébriques finies sur des polynômes
et des fractions. Le théorème des fonctions formelles et une résolution
libre complète ne sont pas requis par ce parcours.
Les résultats de cffr peuvent soutenir les suites exactes, tandis que
HST reste utile pour des généralisations. Pour ce cas précis, la
construction des bases de Gröbner de \cite{CP} et les contractions
sur les trois cartes sont les principaux développements nouveaux.

\begin{thebibliography}{9}
\bibitem{CP} T.~Coquand et H.~Persson, \emph{Gröbner Bases in Type Theory},
TYPES'98, LNCS 1657 (1999), 33--46, section 2, en particulier théorème 11.
Fichier local : \texttt{CPersson.pdf}.
\bibitem{proj} The Stacks Project, section 30.8,
\url{https://stacks.math.columbia.edu/tag/01XS}.
\end{thebibliography}
\end{document}
